\documentclass[11pt,a4paper]{article}
\usepackage[margin=19mm,top=17mm,bottom=17mm]{geometry}
\usepackage{amsmath,amssymb,amsfonts}
\usepackage{mathptmx}
\usepackage{enumitem}
\usepackage{fancyhdr}
\usepackage{array,tabularx,booktabs}
\usepackage[hidelinks]{hyperref}
\setlength{\parindent}{0pt}\setlength{\parskip}{4pt}
\setlist[enumerate,1]{leftmargin=1.1em,itemsep=3pt,parsep=2pt}
\setlist[itemize]{leftmargin=1.4em,itemsep=1pt}
\pagestyle{fancy}\fancyhf{}
\fancyhead[C]{\footnotesize \textbf{Kinematics test - mixed topic}}
\fancyfoot[C]{\footnotesize Page \thepage}
\renewcommand{\headrulewidth}{0.4pt}
\begin{document}
{\centering \large \textbf{Bodha Academy}\par}
{\centering \small \textit{Sector V, Kolkata}\par}
\vspace{2mm}
\noindent \textbf{Marks : 138} \hfill \textbf{Time : 2 : 00 hours}
\noindent\rule{\linewidth}{1.1pt}\vspace{-1mm}\noindent\rule{\linewidth}{0.5pt}
{\centering \Large \textbf{KINEMATICS TEST - MIXED TOPIC --- ANSWER KEY}\par}
{\centering JEE $|$ Physics\par}
\noindent\rule{\linewidth}{1.1pt}\vspace{-1mm}\noindent\rule{\linewidth}{0.5pt}
\vspace{1mm}
\noindent \textbf{Duration:} 2 : 00 hours \quad \textbf{Total Marks:} 138 \quad \textbf{Questions:} 40\par
\vspace{2mm}

\section*{Answer Key}
\noindent\begin{tabularx}{\linewidth}{|c|c|c|X|}\hline
\textbf{Q} & \textbf{Answer} & \textbf{Marks} & \textbf{Concept} \\\hline
1 & $3\,\text{s}$ & 1 & Equations of uniformly accelerated motion \\\hline
2 & $\dfrac{5}{2}$ & 4 & Equations of uniformly accelerated motion \\\hline
3 & $12\,\text{m}$ & 4 & Equations of uniformly accelerated motion \\\hline
4 & $4\,\text{s}$ & 4 & Equations of uniformly accelerated motion \\\hline
5 & $\left(3+\frac{3\sqrt{10}}{2}\right)\,\mathrm{s}$ & 4 & Equations of uniformly accelerated motion \\\hline
6 & 8 s & 4 & Equations of uniformly accelerated motion \\\hline
7 & 3\,\text{m/s} & 4 & Equations of uniformly accelerated motion \\\hline
8 & $4.0\,\mathrm{s}$ & 3 & Equations of uniformly accelerated motion \\\hline
9 & $4\,\mathrm{s}$ & 4 & Equations of uniformly accelerated motion \\\hline
10 & $\dfrac{28}{3}\,\mathrm{s}$ & 4 & Equations of uniformly accelerated motion \\\hline
11 & $-2\,\text{m}$ & 4 & Motion graphs \\\hline
12 & $2\,\text{m}$ & 4 & Motion graphs \\\hline
13 & $12\,\text{m},\ 16\,\text{m}$ & 4 & Motion graphs \\\hline
14 & $-2\,\text{m},\ \dfrac{38}{3}\,\text{m}$ & 4 & Motion graphs \\\hline
15 & $14\ \text{m},\ \dfrac{46}{3}\ \text{m}$ & 4 & Motion graphs \\\hline
16 & $10\,\text{m}$ & 1 & Motion graphs \\\hline
17 & $12\,\text{m},\ 16\,\text{m}$ & 4 & Motion graphs \\\hline
18 & $14\,\text{m},\ \dfrac{46}{3}\,\text{m}$ & 4 & Motion graphs \\\hline
19 & 22 m and 23 m & 1 & Motion graphs \\\hline
20 & $12\,\text{m},\ 24\,\text{m}$ & 4 & Motion graphs \\\hline
21 & $\frac{28}{3}\ \text{m}$ & 4 & Position, velocity and acceleration \\\hline
22 & $\dfrac{13}{5}$ & 1 & Position, velocity and acceleration \\\hline
23 & At $t=5\,\mathrm{s}$, $x=14\,\mathrm{m}$, the velocity is $-2\,\mathrm{m\,s^{-1}}$, and the particle has reversed its direction exactly once. & 1 & Position, velocity and acceleration in one-dimensional motion \\\hline
24 & $\dfrac{28}{3}\ \mathrm{m}$ & 4 & Position, velocity and acceleration \\\hline
25 & $75\,\mathrm{m},\ 2\,\mathrm{s}$ & 4 & Position, velocity and acceleration in one-dimensional motion \\\hline
26 & $-6\,\text{m s}^{-2}$ & 4 & Position, velocity and acceleration \\\hline
27 & \(10.5\,\text{m}\) & 4 & Position, velocity and acceleration \\\hline
28 & $\dfrac{255}{8}\,\mathrm{m}$ & 4 & Position, velocity and acceleration in one-dimensional motion \\\hline
29 & $3:1$ & 4 & Position, velocity and acceleration in one-dimensional motion \\\hline
30 & $180\,\text{m}$ & 4 & Position, velocity and acceleration \\\hline
31 & \(50\,\text{m}\) & 4 & Vertical motion under gravity \\\hline
32 & $3.0\,\text{s}$ & 4 & Vertical motion under gravity \\\hline
33 & $0.6\,\mathrm{s}$ & 1 & Vertical motion under gravity \\\hline
34 & $22.5\,\text{m}$ & 4 & Vertical motion under gravity \\\hline
35 & $0.63\,\mathrm{s}$ & 4 & Vertical motion under gravity \\\hline
36 & It never returns to the release height for any positive time & 1 & Vertical motion under gravity \\\hline
37 & \(7.2\,\mathrm{m}\), with the object below the elevator & 4 & Vertical motion under gravity \\\hline
38 & $\frac{2}{3}\,\text{s}$ & 4 & Vertical motion under gravity \\\hline
39 & $2\,\text{s}$ & 4 & Vertical motion under gravity \\\hline
40 & \(40\,\text{m}\) & 4 & Vertical motion under gravity \\\hline
\end{tabularx}

\section*{Solutions}
\begin{enumerate}
\item The positions at time $t$ are $x_A=0+2t+\frac12(1)t^2=2t+\frac12t^2$ and $x_B=27-4t+\frac12(-1)t^2=27-4t-\frac12t^2$. At the meeting point, $x_A=x_B$, so $2t+\frac12t^2=27-4t-\frac12t^2$, giving $t^2+6t-27=0$. Hence $(t-3)(t+9)=0$, so $t=3\,\text{s}$ or $t=-9\,\text{s}$. Since time must be positive, the meeting occurs at $t=3\,\text{s}$. Also, the separation is $x_B-x_A=27-6t-t^2$, whose derivative is $-6-2t<0$ for $t\geq0$, so it decreases continuously from $27\,\text{m}$ and can reach zero only once for positive time. Thus this is the first meeting time.
\item For uniformly accelerated motion, the average velocity over an interval equals the instantaneous velocity at the midpoint of that interval. The midpoints of the three equal intervals are separated by equal time intervals, so the corresponding average velocities form an arithmetic progression. Hence, $\bar v_3=2\bar v_2-\bar v_1=2(14)-8=20\,\mathrm{m\,s^{-1}}$. Therefore, $s_1=\bar v_1T=8T$ and $s_3=\bar v_3T=20T$. Thus, $\dfrac{s_3}{s_1}=\dfrac{20T}{8T}=\dfrac{5}{2}$, so the correct option is $\dfrac{5}{2}$.
\item After time $t$, the displacement of A is $s_A=\frac{1}{2}(2)t^2=t^2$, and the displacement of B toward A is $s_B=\frac{1}{2}(1)t^2=\frac{t^2}{2}$. At the meeting point, their total displacement equals the initial separation: $t^2+\frac{t^2}{2}=18$. Thus, $\frac{3}{2}t^2=18$, giving $t^2=12$. Therefore, the meeting position measured from A's initial position is $x=s_A=t^2=12\,\text{m}$.
\item Let $t$ be the time elapsed after particle B starts. At this instant, particle A has been moving for $(t+2)\,\text{s}$. Hence, measured from the common marker, their positions are
\[
x_A=6(t+2)+\frac12(2)(t+2)^2,
\]
\[
x_B=10t+\frac12(4)t^2.
\]
On equating positions,
\[
6(t+2)+(t+2)^2=10t+2t^2.
\]
Expanding and simplifying,
\[
6t+12+t^2+4t+4=10t+2t^2,
\]
\[
t^2=16.
\]
Thus $t=\pm4\,\text{s}$. Since $t$ is measured after particle B starts, the physically admissible value is $t=4\,\text{s}$. Therefore, they meet for the first time $4\,\text{s}$ after B starts.
\item Let $t$ be the time elapsed after cyclist A passes the marker. The positions measured from the marker are $x_A=6t+\frac{1}{2}(2)t^2=6t+t^2$ and, for $t\geq 3\,\mathrm{s}$, $x_B=\frac{1}{2}(6)(t-3)^2=3(t-3)^2$. At the catching instant, $x_A=x_B$, so $6t+t^2=3(t-3)^2$. This gives $2t^2-24t+27=0$, whose roots are $t=6\pm\frac{3\sqrt{10}}{2}\,\mathrm{s}$. The smaller root is approximately $1.26\,\mathrm{s}$ and is invalid because cyclist B has not yet started. Hence $t=6+\frac{3\sqrt{10}}{2}\,\mathrm{s}$. Therefore, the time after B starts is $t-3=\left(3+\frac{3\sqrt{10}}{2}\right)\,\mathrm{s}$.
\item Choose the initial position of A as x = 0. Then, the positions after time t are x_A = 10t + (1/2)(5)t^2 = 10t + 2.5t^2 and x_B = 24 + 15t + (1/2)(3)t^2 = 24 + 15t + 1.5t^2. At the instant of catching, x_A = x_B. Hence, 10t + 2.5t^2 = 24 + 15t + 1.5t^2, giving t^2 - 5t - 24 = 0. Therefore, (t - 8)(t + 3) = 0, so t = 8 s or t = -3 s. Since time cannot be negative, the first catching time is 8 s.
\item For uniformly accelerated motion, x(t)=ut+\frac{1}{2}at^2. Since x=0 at t=0 and t=6\,\text{s}, x(t) can be written as x(t)=K t(t-6). Using x(2)=12\,\text{m}, we get 12=K(2)(-4), so K=-\frac{3}{2}. Therefore, x(t)=9t-\frac{3}{2}t^2. To find the times when x=12\,\text{m}, solve 9t-\frac{3}{2}t^2=12, giving t^2-6t+8=0 and hence t=2\,\text{s} or 4\,\text{s}. The second passage occurs at t=4\,\text{s}. The velocity is v=\frac{dx}{dt}=9-3t, so v(4)=-3\,\text{m/s}. Thus, the speed is |v|=3\,\text{m/s}.
\item The positions of A and B at time $t$ are, respectively, $x_A=4t+\frac{1}{2}(2)t^2=4t+t^2$ and $x_B=56-\frac{1}{2}(3)t^2=56-\frac{3}{2}t^2$. At interception, $x_A=x_B$, so $4t+t^2=56-\frac{3}{2}t^2$, giving $\frac{5}{2}t^2+4t-56=0$. Thus $5t^2+8t-112=0$, whose roots are $t=4\,\mathrm{s}$ and $t=-\frac{28}{5}\,\mathrm{s}$. The negative root is physically inadmissible, so the particles meet after $\boxed{4.0\,\mathrm{s}}$.
\item For Cart A, its position at time $t$ is $x_A=10t+\frac{1}{2}(4)t^2=10t+2t^2$. For $t>3\,\mathrm{s}$, Cart B has been moving for $(t-3)$ seconds, so its position is $x_B=71+\frac{1}{2}(2)(t-3)^2=71+(t-3)^2$. At the catching time, $x_A=x_B$: $10t+2t^2=71+(t-3)^2$. On simplifying, $t^2+16t-80=0$, giving $(t-4)(t+20)=0$. The roots are $t=4\,\mathrm{s}$ and $t=-20\,\mathrm{s}$. The negative root is unphysical, while $t=4\,\mathrm{s}>3\,\mathrm{s}$ satisfies the condition that Cart B has already been released. Hence, Cart A catches Cart B at $4\,\mathrm{s}$.
\item Using $v^2=u^2+2as$ with $v=0$, $u=18\,\mathrm{m\,s^{-1}}$, and $s=120\,\mathrm{m}$: $0=18^2+2a(120)$, giving $a=-1.35\,\mathrm{m\,s^{-2}}$. The total time taken to stop after entering the platform is obtained from $v=u+at$: $0=18-1.35t$, so $t=\dfrac{40}{3}\,\mathrm{s}$. Since the warning signal is received $4\,\mathrm{s}$ after entry, the remaining time is $\dfrac{40}{3}-4=\dfrac{28}{3}\,\mathrm{s}$.
\item The displacement equals the signed area under the velocity–time graph. For $0\le t\le2\,\text{s}$, the area is $\frac12(2)(4)=4\,\text{m}$. For $2\le t\le4\,\text{s}$, the area is $-\frac12(2)(2)=-2\,\text{m}$. For $4\le t\le6\,\text{s}$, the area is $(-2)(2)=-4\,\text{m}$. Hence, $\Delta x=4-2-4=-2\,\text{m}$)
\item Displacement equals the signed area under the velocity-time graph. During $0$ to $2\,\text{s}$, the area is $\frac{1}{2}\times 2\times 4=4\,\text{m}$. During $2$ to $4\,\text{s}$, the velocity changes linearly from $4$ to $-2\,\text{m s}^{-1}$, so the signed trapezoidal area is $\frac{1}{2}(4-2)\times 2=2\,\text{m}$. During $4$ to $6\,\text{s}$, the area is $(-2)\times 2=-4\,\text{m}$. Hence, the net displacement is $4+2-4=2\,\text{m}$, so the correct option is $\boxed{2\,\text{m}}$.
\item The signed area under the velocity–time graph gives displacement. For $0\le t\le2$, the area is $\frac12(2)(4)=4\,\text{m}$. For $2\le t\le4$, it is $(2)(4)=8\,\text{m}$. During the last $2\,\text{s}$, velocity changes linearly from $4$ to $-4$ and becomes zero halfway through, so the positive and negative triangular areas are equal and cancel. Hence, displacement $=4+8=12\,\text{m}$. For distance, all areas are taken positive: $4+8+\frac12(1)(4)+\frac12(1)(4)=16\,\text{m}$.
\item The displacement is the signed area under the velocity–time graph. For $0\leq t\leq4$, the velocity changes from $4$ to $-2\,\text{m s}^{-1}$, so the net displacement is $\left(\frac{4+(-2)}{2}\right)\times4=4\,\text{m}$. From $t=4$ to $6$, the displacement is $(-2)(2)=-4\,\text{m}$, and from $t=6$ to $8$, it is $\left(\frac{-2+0}{2}\right)\times2=-2\,\text{m}$. Hence, $s=4-4-2=-2\,\text{m}$. The velocity becomes zero at $t=\frac{8}{3}\,\text{s}$ during the first interval. The positive triangular area is $\frac12\times\frac83\times4=\frac{16}{3}\,\text{m}$, while the negative triangular area from $t=\frac83$ to $4$ is $\frac12\times\frac43\times2=\frac43\,\text{m}$. Adding the magnitudes of all portions, the total distance is $\frac{16}{3}+\frac{4}{3}+4+2=\frac{38}{3}\,\text{m}$.
\item Displacement equals the algebraic area under the velocity–time graph. From $0$ to $2$ s, the area is $\frac12(2)(4)=4$ m. From $2$ to $4$ s, it is $(2)(4)=8$ m. From $4$ to $6$ s, the signed trapezoidal area is $\frac{(4+(-2))}{2}\times 2=2$ m. Hence, displacement $=4+8+2=14$ m.\n\nFor distance, the velocity becomes zero on the last segment. Since the velocity decreases from $4$ to $-2$ in $2$ s, it becomes zero after $\frac{4}{3}$ s. The positive triangular area in this segment is $\frac12\times\frac43\times4=\frac83$ m, and the magnitude of the negative triangular area is $\frac12\times\frac23\times2=\frac23$ m. Therefore, total distance $=4+8+\frac83+\frac23=\frac{46}{3}$ m.
\item Displacement equals the algebraic area under the velocity–time graph. From $0$ to $2\,\text{s}$, the area is $\frac{1}{2}(2)(4)=4\,\text{m}$. From $2$ to $4\,\text{s}$, it is $(2)(4)=8\,\text{m}$. From $4$ to $6\,\text{s}$, the velocity changes linearly from $4$ to $-2\,\text{m s}^{-1}$, so the signed area is $\frac{(4+(-2))}{2}\times 2=2\,\text{m}$. From $6$ to $8\,\text{s}$, the displacement is $(-2)(2)=-4\,\text{m}$. Hence, total displacement $=4+8+2-4=10\,\text{m}$.
\item Displacement equals the algebraic area under the velocity–time graph. During $0$–$2\,\text{s}$, the area is $\frac{1}{2}\times 2\times 4=4\,\text{m}$. During $2$–$4\,\text{s}$, it is $2\times 4=8\,\text{m}$. From $4$–$6\,\text{s}$, velocity changes linearly from $4$ to $-4\,\text{m s}^{-1}$, so the net signed area is zero. Hence, displacement $=4+8=12\,\text{m}$. For distance, the final interval consists of two triangles, each with base $1\,\text{s}$ and height $4\,\text{m s}^{-1}$, contributing $4\,\text{m}$ in total. Therefore, total distance $=4+8+4=16\,\text{m}$.
\item Displacement equals the algebraic area under the velocity–time graph. From $0$ to $2\,\text{s}$, the area is $\frac{1}{2}\times2\times4=4\,\text{m}$. From $2$ to $4\,\text{s}$, it is $2\times4=8\,\text{m}$. From $4$ to $6\,\text{s}$, the signed trapezium area is $\frac{1}{2}(4+(-2))\times2=2\,\text{m}$. Hence, displacement $=4+8+2=14\,\text{m}$. For distance, the velocity becomes zero during the last segment. Since it falls from $4$ to $-2$ in $2\,\text{s}$, the zero-velocity instant is $\frac{4}{6}\times2=\frac{4}{3}\,\text{s}$ after $t=4\,\text{s}$. The positive and negative triangular areas in this segment are $\frac{1}{2}\times\frac{4}{3}\times4=\frac{8}{3}\,\text{m}$ and $\frac{1}{2}\times\frac{2}{3}\times2=\frac{2}{3}\,\text{m}$, respectively. Therefore, distance $=4+8+\frac{8}{3}+\frac{2}{3}=\frac{46}{3}\,\text{m}$.
\item The displacement equals the algebraic area under the velocity–time graph. For 0–2 s, the area is \(\frac{1}{2}\times 2\times 6=6\,\text{m}\). For 2–4 s, it is \(2\times 6=12\,\text{m}\). During 4–6 s, the velocity falls from 6 m/s to -2 m/s, so the net area is \(\frac{1}{2}(6+(-2))\times 2=4\,\text{m}\). Hence, displacement \(=6+12+4=22\,\text{m}\). The velocity becomes zero after \(\frac{6}{4}=1.5\,\text{s}\), i.e. at \(t=5.5\,\text{s}\). Therefore, the positive area from 4–5.5 s is \(\frac{1}{2}\times 1.5\times 6=4.5\,\text{m}\), and the magnitude of the negative area from 5.5–6 s is \(\frac{1}{2}\times 0.5\times 2=0.5\,\text{m}\). Thus, total distance \(=6+12+4.5+0.5=23\,\text{m}\).
\item The displacement equals the algebraic area under the velocity–time graph. From $0$ to $2$ s, the area is $\frac12(2)(4)=4$ m. From $2$ to $5$ s, it is $(3)(4)=12$ m. From $5$ to $7$ s, the velocity changes linearly from $4$ to $-4\,\text{m s}^{-1}$, so the positive and negative triangular areas cancel, giving zero net displacement. From $7$ to $9$ s, the displacement is $-\frac12(2)(4)=-4$ m. Hence, $s=4+12+0-4=12$ m. For distance, areas are added without signs: $4+12+\frac12(2)(4)+\frac12(2)(4)=24$ m. Therefore, the required pair is $12$ m and $24$ m.
\item The velocity becomes zero when $t^2-4t+3=0$, i.e. $(t-1)(t-3)=0$. Thus, the particle may change direction at $t=1\ \text{s}$ and $t=3\ \text{s}$. Since $v(t)>0$ for $0\le t<1$, $v(t)<0$ for $1<t<3$, and $v(t)>0$ for $3<t\le5$, calculate the displacement in each interval. An antiderivative of $v(t)$ is $F(t)=\frac{t^3}{3}-2t^2+3t$. Hence, $\Delta x_{0\to1}=F(1)-F(0)=\frac{4}{3}\ \text{m}$, $\Delta x_{1\to3}=F(3)-F(1)=-\frac{4}{3}\ \text{m}$, and $\Delta x_{3\to5}=F(5)-F(3)=\frac{20}{3}\ \text{m}$. Therefore, total distance $=\left|\frac{4}{3}\right|+\left|-\frac{4}{3}\right|+\left|\frac{20}{3}\right|=\frac{28}{3}\ \text{m}$.
\item The direction reverses when $v(t)=0$: $6-2t=0$, giving $t=3\,\mathrm{s}$. The displacement from $0$ to $3\,\mathrm{s}$ is $\Delta x_1=\int_0^3(6-2t)\,dt=[6t-t^2]_0^3=9\,\mathrm{m}$. The displacement from $3$ to $5\,\mathrm{s}$ is $\Delta x_2=\int_3^5(6-2t)\,dt=[6t-t^2]_3^5=5-9=-4\,\mathrm{m}$. Hence, total distance $=|9|+|-4|=13\,\mathrm{m}$, while net displacement $=9-4=5\,\mathrm{m}$. Therefore, the required ratio is $\dfrac{13}{5}$.
\item From $t=0$ to $t=3\,\mathrm{s}$, the displacement is $\Delta x_1=(4)(3)=12\,\mathrm{m}$. During the interval $3\,\mathrm{s}\le t\le5\,\mathrm{s}$, the velocity varies uniformly, so the average velocity is $\bar v=\frac{4+(-2)}{2}=1\,\mathrm{m\,s^{-1}}$. Hence, $\Delta x_2=(1)(2)=2\,\mathrm{m}$. Therefore, $x(5)=12+2=14\,\mathrm{m}$. Since the velocity changes continuously from positive to negative, it becomes zero once and the particle reverses direction exactly once before $t=5\,\mathrm{s}$. At $t=5\,\mathrm{s}$, its velocity is $-2\,\mathrm{m\,s^{-1}}$.
\item Factor the velocity: $v(t)=(t-1)(t-3)$. Thus, the particle moves in the positive direction on $[0,1]$, in the negative direction on $[1,3]$, and again in the positive direction on $[3,5]$.\n\nThe position is obtained by integration:\n\n$x(t)=\int_0^t (u^2-4u+3)\,du=\dfrac{t^3}{3}-2t^2+3t$.\n\nHence, $x(0)=0$, $x(1)=\dfrac{4}{3}\ \mathrm{m}$, $x(3)=0$, and $x(5)=\dfrac{20}{3}\ \mathrm{m}$. Therefore, the total distance travelled is\n\n$|x(1)-x(0)|+|x(3)-x(1)|+|x(5)-x(3)|=\dfrac{4}{3}+\dfrac{4}{3}+\dfrac{20}{3}=\dfrac{28}{3}\ \mathrm{m}$.
\item In the first stage, $u=0$, $a=2\,\mathrm{m\,s^{-2}}$, and $t_1=5\,\mathrm{s}$. Hence, the velocity attained is $v=u+at_1=0+2(5)=10\,\mathrm{m\,s^{-1}}$, and the displacement is $s_1=ut_1+\frac{1}{2}at_1^2=\frac{1}{2}(2)(5^2)=25\,\mathrm{m}$. In the second stage, the particle moves with constant velocity $10\,\mathrm{m\,s^{-1}}$ for $4\,\mathrm{s}$, so $s_2=10\times4=40\,\mathrm{m}$. During the final stage, the initial velocity is $10\,\mathrm{m\,s^{-1}}$ and the acceleration is $-5\,\mathrm{m\,s^{-2}}$. From $v=u+at_3$, $0=10-5t_3$, giving $t_3=2\,\mathrm{s}$. Its displacement is $s_3=\frac{u+v}{2}t_3=\frac{10+0}{2}\times2=10\,\mathrm{m}$. Therefore, the total displacement is $s_1+s_2+s_3=25+40+10=75\,\mathrm{m}$, and the duration of the final stage is $2\,\mathrm{s}$.
\item Let the acceleration during $0\le t\le2\,\text{s}$ be $a_1$. Since $u=0$ and the displacement in this interval is $12\,\text{m}$, $12=\frac12a_1(2)^2$, giving $a_1=6\,\text{m s}^{-2}$. Hence, the velocity at $t=2\,\text{s}$ is $v_2=u+a_1t=12\,\text{m s}^{-1}$. Velocity is continuous even though acceleration changes abruptly. For the second interval, $0=v_2+a_2(2)$ because the final velocity at $t=4\,\text{s}$ is zero. Thus, $0=12+2a_2$, so $a_2=-6\,\text{m s}^{-2}$. In this interval, $v=12-6(t-2)>0$ for $2<t<4$, so the stated motion is consistent. Also, the displacement in the second interval is $\frac{12+0}{2}\times2=12\,\text{m}$, giving the stated total displacement of $24\,\text{m}$.
\item During the first phase, \(v=u+a t\). The particle stops when \(0=6-3t_1\), giving \(t_1=2\,\text{s}\). Its displacement in this phase is \(s_1=ut_1+\frac12at_1^2=6(2)+\frac12(-3)(2)^2=6\,\text{m}\). The remaining time is \(5-2=3\,\text{s}\). In the second phase, the particle starts from rest with acceleration \(1\,\text{m s}^{-2}\), so \(s_2=\frac12(1)(3)^2=4.5\,\text{m}\). Therefore, its position at \(t=5\,\text{s}\) is \(x=s_1+s_2=6+4.5=10.5\,\text{m}\).
\item During the first phase, $u=0$, $a=2\,\mathrm{m\,s^{-2}}$ and $t=5\,\mathrm{s}$. Hence, the speed attained is $v=u+at=10\,\mathrm{m\,s^{-1}}$, which is the maximum speed in the motion. The displacement in this phase is $s_1=ut+\frac12at^2=\frac12(2)(5^2)=25\,\mathrm{m}$.\

During deceleration, the particle comes to rest from $10\,\mathrm{m\,s^{-1}}$ with acceleration $-5\,\mathrm{m\,s^{-2}}$. Its displacement is $s_2=\frac{v^2-u^2}{2a}=\frac{0-100}{2(-5)}=10\,\mathrm{m}$. Thus, its position at the instant of reversal is $x_0=25+10=35\,\mathrm{m}$.\

One-fourth of the maximum speed is $10/4=2.5\,\mathrm{m\,s^{-1}}$. During reversal, the acceleration is $-1\,\mathrm{m\,s^{-2}}$, starting from rest. The time required to attain speed $2.5\,\mathrm{m\,s^{-1}}$ is $t=2.5\,\mathrm{s}$. The displacement during this reversed segment is $s_3=\frac12(-1)(2.5)^2=-\frac{25}{8}\,\mathrm{m}$. Therefore, the displacement from the starting point is $x=x_0+s_3=35-\frac{25}{8}=\frac{255}{8}\,\mathrm{m}$.
\item The velocity is $v(t)=\frac{dx}{dt}=3t^2-12t+9=3(t-1)(t-3)$. Hence, the particle reverses direction at $t=1$ s and $t=3$ s. Its positions at the relevant instants are $x(0)=2$ m, $x(1)=6$ m, $x(3)=2$ m, and $x(4)=6$ m. Therefore, the total distance travelled is $|6-2|+|2-6|+|6-2|=4+4+4=12$ m. The net displacement is $x(4)-x(0)=6-2=4$ m, whose magnitude is $4$ m. Thus, the required ratio is $\frac{12}{4}=3:1$.
\item The position functions are $x_A=8t+\frac12(2)t^2=8t+t^2$ and $x_B=30+20t+\frac12(-1)t^2=30+20t-\frac12t^2$. At the meeting point, $x_A=x_B$, so $8t+t^2=30+20t-\frac12t^2$, or $3t^2-24t-60=0$. Thus, $t=10\,\text{s}$ or $t=-2\,\text{s}$. The physically relevant time after $t=0$ is $t=10\,\text{s}$. Therefore, $x=8(10)+(10)^2=180\,\text{m}$.
\item The ball's initial velocity relative to the ground is the sum of its velocity relative to the elevator and the elevator's velocity: \(u=15+5=20\,\text{m s}^{-1}\) upward. At the highest point, its velocity becomes zero. Using \(v^2=u^2-2g\Delta h\), \(0=20^2-2(10)\Delta h\), so \(\Delta h=20\,\text{m}\). Therefore, the maximum height above the ground is \(30+20=50\,\text{m}\).
\item The slope of the velocity--time graph gives the acceleration: $a=\dfrac{\Delta v}{\Delta t}=-g$. The change in velocity between the two instants is $\Delta v=(-12)-(+18)=-30\,\text{m s}^{-1}$. Therefore, $-30=-10\,\Delta t$, giving $\Delta t=3\,\text{s}$. The sign of velocity only identifies the direction of motion; it does not affect the elapsed-time calculation.
\item Take the downward direction relative to the elevator as positive. Initially, the ball has zero velocity relative to the elevator. The ball has downward acceleration $g$, while the elevator has upward acceleration $2\,\mathrm{m\,s^{-2}}$. Therefore, the relative downward acceleration of the ball with respect to the elevator is $a_{\mathrm{rel}}=g+2=12\,\mathrm{m\,s^{-2}}$. Using $s=u t+\frac{1}{2}a_{\mathrm{rel}}t^2$, we get $2.16=0+\frac{1}{2}(12)t^2$, so $t^2=0.36$ and hence $t=0.6\,\mathrm{s}$.
\item The velocity at time $t$ is $v=15-10(t-1)$. The ball comes momentarily to rest when $v=0$: $15-10(t-1)=0$, giving $t=2.5\,\text{s}$. From $t=1$ to $t=2.5\,\text{s}$, the ball moves upward. The distance travelled is the area of the triangular region under the speed–time graph: $s_1=\frac{1}{2}(1.5)(15)=11.25\,\text{m}$. From $t=2.5$ to $t=4\,\text{s}$, it moves downward, with speed increasing from $0$ to $15\,\text{m s}^{-1}$. Thus, $s_2=\frac{1}{2}(1.5)(15)=11.25\,\text{m}$. Therefore, the total distance is $s_1+s_2=11.25+11.25=22.5\,\text{m}$.
\item In the non-inertial frame of the accelerating lift, the coin experiences an effective downward acceleration equal to $g+a_L$, where $a_L=2\,\mathrm{m\,s^{-2}}$ is the upward acceleration of the lift. Thus, $a_{\mathrm{rel}}=10+2=12\,\mathrm{m\,s^{-2}}$. Since the coin is released from rest relative to the lift, its initial relative velocity is zero. For the relative displacement of $2.4\,\mathrm{m}$, $2.4=\frac{1}{2}(12)t^2$, giving $t^2=0.4$ and $t=\sqrt{0.4}\approx 0.63\,\mathrm{s}$.
\item Let upward be positive and let $\tau$ denote the time measured from the instant the elevator starts accelerating. During the initial $2\,\text{s}$, the ball has zero initial velocity relative to the elevator and relative acceleration $-g=-10\,\text{m s}^{-2}$. Hence, at the instant the elevator starts accelerating, its velocity relative to the elevator is $v_{\text{rel},0}=-10(2)=-20\,\text{m s}^{-1}$, and its displacement from the release point is $y_{\text{rel},0}=\frac12(-10)(2)^2=-20\,\text{m}$.\n\nAfterwards, the elevator has upward acceleration $+2\,\text{m s}^{-2}$ while the ball has acceleration $-10\,\text{m s}^{-2}$. Therefore, the ball's relative acceleration is $a_{\text{rel}}=-10-2=-12\,\text{m s}^{-2}$. Its relative displacement is thus\n\n$$y_{\text{rel}}=-20-20\tau-\frac12(12)\tau^2=-20-20\tau-6\tau^2.$$\n\nFor every $\tau>0$, this expression is negative, so the ball continues to move below the release point relative to the elevator. Setting $y_{\text{rel}}=0$ gives $6\tau^2+20\tau+20=0$, whose roots are negative or nonphysical. Hence, there is no positive time at which the ball returns to the release height.
\item At the instant of release, the object has the same upward velocity as the elevator, \(u=8\,\mathrm{m\,s^{-1}}\). Taking the release point as the origin, the object's position after time \(t\) is \(y_o=8t-\frac{1}{2}gt^2\), while the elevator's position is \(y_e=8t\), since it continues moving upward uniformly. Thus, \(y_o-y_e=-\frac{1}{2}gt^2\). For \(t=1.2\,\mathrm{s}\), \(y_o-y_e=-\frac{1}{2}(10)(1.2)^2=-7.2\,\mathrm{m}\). Therefore, the object is \(7.2\,\mathrm{m}\) below the elevator.
\item Choose the elevator as the reference frame. The ball initially has relative velocity $u=4\,\text{m s}^{-1}$ upward. Its acceleration relative to the elevator is $a_{\text{rel}}=-(g+2)=-(10+2)=-12\,\text{m s}^{-2}$, since the ball accelerates downward at $10\,\text{m s}^{-2}$ while the elevator accelerates upward at $2\,\text{m s}^{-2}$. For the ball to return to its initial level relative to the elevator, its relative displacement must be zero: $0=4t-\frac{1}{2}(12)t^2=t(4-6t)$. Excluding the initial instant $t=0$, we obtain $t=\frac{2}{3}\,\text{s}$, which is the required time.
\item In the reference frame of the accelerating lift, the ball experiences an effective downward acceleration equal to $g+a_{\text{lift}}=10+2=12\,\text{m s}^{-2}$. Its initial velocity relative to the passenger is $u_{\text{rel}}=12\,\text{m s}^{-1}$ upward. For the ball to return to its initial relative position, the time is $t=\frac{2u_{\text{rel}}}{g+a_{\text{lift}}}=\frac{2\times12}{12}=2\,\text{s}$.
\item The velocity of the ball relative to the ground immediately after projection is the sum of the lift's velocity and the ball's velocity relative to the lift: \(u=5+15=20\,\text{m s}^{-1}\) upward. At the highest point, the final velocity is zero. Using \(v^2=u^2-2g\Delta h\), \(0=20^2-2(10)\Delta h\), giving \(\Delta h=20\,\text{m}\). Since the ball is released at a height of \(20\,\text{m}\), its maximum height above the ground is \(20+20=40\,\text{m}\).
\end{enumerate}
\end{document}